\documentclass{article}
\usepackage{amsmath,amssymb,amsthm}
\usepackage{enumitem}
\newtheorem{theorem}{Theorem}
\newtheorem{lemma}{Lemma}
\newcommand{\E}{\mathbb{E}}
\newcommand{\R}{\mathbb{R}}
\newcommand{\norm}[1]{\left\lVert#1\right\rVert}
\begin{document}

\section*{Algorithm Name}
\textbf{Stochastic Coordinate Momentum Descent (SCMD)}  

At each iteration a single coordinate is drawn according to a fixed
distribution  \(p=(p_{1},\dots ,p_{d})\)  (\(p_{i}>0,\;\sum_{i}p_{i}=1\)).
Only that coordinate is updated, and a heavy‑ball momentum term is kept
for every coordinate.

-----------------------------------------------------------------------

\section*{Mathematical Formulation}
Let \(f:\R^{d}\rightarrow\R\) be the objective.  
For each coordinate \(i\) we keep a momentum variable
\(m_{i}^{k}\in\R\).  The algorithm proceeds as follows.

\begin{itemize}
  \item \textbf{Sampling:}  draw \(i_{k}\in\{1,\dots ,d\}\) with
        \(\Pr(i_{k}=i)=p_{i}\).
  \item \textbf{Gradient on the sampled coordinate:}
        \[
        g_{i_{k}}^{k}= \frac{\partial f}{\partial w_{i_{k}}}\bigl(w^{k}\bigr).
        \]
  \item \textbf{Momentum update (heavy‑ball):}
        \[
        m_{i_{k}}^{k+1}= \beta\, m_{i_{k}}^{k}+(1-\beta)\, g_{i_{k}}^{k},
        \qquad \beta\in[0,1).
        \]
        For all \(j\neq i_{k}\) set \(m_{j}^{k+1}=m_{j}^{k}\).
  \item \textbf{Coordinate step:}
        \[
        w_{i_{k}}^{k+1}= w_{i_{k}}^{k}-\eta_{i_{k}}\, m_{i_{k}}^{k+1},
        \qquad 
        \eta_{i}= \frac{1}{L_{i}(1-\beta)} ,
        \]
        where \(L_{i}>0\) is the \emph{coordinate‑wise smoothness}
        constant (see Assumption 2).  All other coordinates stay unchanged,
        i.e. \(w_{j}^{k+1}=w_{j}^{k}\) for \(j\neq i_{k}\).
\end{itemize}

The full iterate vector after iteration \(k\) is denoted
\(w^{k}=(w_{1}^{k},\dots ,w_{d}^{k})\).

-----------------------------------------------------------------------

\section*{Pseudocode}
\begin{center}
\begin{tabular}{l}
\hline
\textbf{Input:}  initial point \(w^{0}\in\R^{d}\), momentum \(\beta\in[0,1)\),\\
\qquad\qquad   coordinate smoothness constants \(\{L_{i}\}_{i=1}^{d}\),\\
\qquad\qquad   probability vector \(p\) with \(p_{i}>0,\;\sum_{i}p_{i}=1\).\\
\textbf{Initialize:}  \(m^{0}=0\in\R^{d}\).\\
\textbf{for} \(k=0,1,2,\dots ,K-1\) \textbf{do}\\
\qquad Sample \(i_{k}\sim p\).\\
\qquad Compute \(g_{i_{k}}^{k}= \partial_{i_{k}}f(w^{k})\).\\
\qquad Update momentum \(\displaystyle 
        m_{i_{k}}^{k+1}= \beta m_{i_{k}}^{k}+(1-\beta)g_{i_{k}}^{k}\).\\
\qquad Set step size \(\displaystyle \eta_{i_{k}}=
        \frac{1}{L_{i_{k}}(1-\beta)}\).\\
\qquad Update coordinate 
        \(\displaystyle w_{i_{k}}^{k+1}= w_{i_{k}}^{k}-\eta_{i_{k}}m_{i_{k}}^{k+1}\).\\
\qquad For \(j\neq i_{k}\): \(w_{j}^{k+1}=w_{j}^{k},\;m_{j}^{k+1}=m_{j}^{k}\).\\
\textbf{end for}\\
\textbf{Output:} \(w^{K}\).\\
\hline
\end{tabular}
\end{center}

-----------------------------------------------------------------------

\section*{Dry Run Example}
Consider the one‑dimensional problem  
\(f(w)=(w-3)^{2}\) (so \(d=1\)).  
We have \(L_{1}=2\) (since \(f''(w)=2\)), \(\mu=2\) (strong convexity),  
\(p_{1}=1\), \(\beta=0\) (no momentum for illustration) and
\(\eta_{1}=1/(L_{1}(1-\beta))=1/2\).

\begin{center}
\begin{tabular}{c|c|c|c}
\textbf{Iter.} & \(w^{k}\) & Gradient \(g^{k}=2(w^{k}-3)\) & \(f(w^{k})\)\\\hline
0 & 0 & \(-6\) & 9\\
1 & \(0-\frac12(-6)=3\) & \(0\) & 0\\
2 & \(3-\frac12(0)=3\) & \(0\) & 0\\
3 & \(3-\frac12(0)=3\) & \(0\) & 0\\
\end{tabular}
\end{center}
After the first iteration the optimum \(w^{*}=3\) is reached and the
objective stays at 0.

-----------------------------------------------------------------------

\section*{Assumptions}
\begin{enumerate}[label=\textbf{A\arabic*:},leftmargin=*]
\item \textbf{Strong convexity.}  
\(f\) is \(\mu\)-strongly convex:
\[
f(y)\ge f(x)+\langle\nabla f(x),y-x\rangle+\frac{\mu}{2}\|y-x\|^{2},
\qquad \forall x,y\in\R^{d},
\]
with \(\mu>0\).

\item \textbf{Coordinate‑wise smoothness.}  
For each \(i\) there exists \(L_{i}>0\) such that for all
\(x\in\R^{d}\) and any scalar \(\delta\):
\[
\bigl|\partial_{i}f(x+\delta e_{i})-\partial_{i}f(x)\bigr|
   \le L_{i}\,|\delta|,
\]
where \(e_{i}\) is the \(i\)-th unit vector.  Equivalently,
\[
f(x+\delta e_{i})\le f(x)+\partial_{i}f(x)\,\delta
      +\frac{L_{i}}{2}\,\delta^{2}.
\]

\item \textbf{Sampling distribution.}  
\(p_{i}>0\) for all \(i\) and \(\sum_{i=1}^{d}p_{i}=1\).

\item \textbf{Momentum parameter.}  
\(\beta\in[0,1)\) is fixed and the stepsizes are chosen as
\(\displaystyle \eta_{i}= \frac{1}{L_{i}(1-\beta)}\).
\end{enumerate}

-----------------------------------------------------------------------

\section*{Main Convergence Theorem}
\begin{theorem}[Linear convergence of SCMD]\label{thm:linear}
Under Assumptions A1–A4, let \(\{w^{k}\}_{k\ge0}\) be the sequence
generated by SCMD.  Define the weighted condition number
\[
\kappa_{\max}\;:=\;\max_{i}\frac{L_{i}}{\mu},
\qquad
\rho\;:=\; \min_{i}\;p_{i}\,\frac{1-\beta}{\kappa_{\max}} \;\in\;(0,1).
\]
Then for all \(k\ge0\)
\[
\E\!\left[f(w^{k})-f(w^{*})\right]\;\le\;
(1-\rho)^{k}\,\bigl(f(w^{0})-f(w^{*})\bigr).
\]
Consequently,
\(\E\!\bigl[\|w^{k}-w^{*}\|^{2}\bigr]\le
\frac{2}{\mu}(1-\rho)^{k}\bigl(f(w^{0})-f(w^{*})\bigr)\),
i.e. SCMD enjoys a \emph{geometric} (linear) rate in expectation.
\end{theorem}

-----------------------------------------------------------------------

\section*{Theoretical Properties}
\begin{itemize}
\item \textbf{Stability.}  The stepsize choice \(\eta_{i}=1/(L_{i}(1-\beta))\)
guarantees that the one‑step surrogate (coordinate smoothness inequality)
is always an upper bound, preventing divergence irrespective of the
sampling probabilities.

\item \textbf{Hyper‑parameter sensitivity.}
  \begin{itemize}
  \item Larger momentum \(\beta\) enlarges the effective stepsize
    \(1/(1-\beta)\) but simultaneously reduces the factor \((1-\beta)\)
    in \(\rho\); the theorem shows that any \(\beta<1\) preserves linear
    convergence, yet the optimal rate is attained for a moderate
    \(\beta\) that balances the two effects.
  \item Non‑uniform probabilities: a coordinate with a larger
    probability \(p_{i}\) contributes more to \(\rho\); in practice,
    setting \(p_{i}\propto L_{i}^{-1}\) (importance sampling) maximises
    \(\rho\) and yields the fastest theoretical rate.
  \end{itemize}

\item \textbf{Comparison to baselines.}
  \begin{itemize}
  \item If \(\beta=0\) the method reduces to standard stochastic
    coordinate descent; the rate in Theorem~\ref{thm:linear}
    coincides with the classical \(\bigl(1-\frac{\mu}{\sum_{i}L_{i}}\bigr)^{k}\)
    bound.
  \item With uniform sampling (\(p_{i}=1/d\)) and identical smoothness
    (\(L_{i}=L\)) the rate becomes \(\bigl(1-\frac{1-\beta}{d\kappa}\bigr)^{k}\),
    a speed‑up of factor \(1/(1-\beta)\) over the non‑momentum version.
  \end{itemize}
\end{itemize}

-----------------------------------------------------------------------

\section*{Preliminaries (Definitions and Lemmas)}
\begin{lemma}[One‑step progress for a fixed coordinate]\label{lem:one-step}
Let \(i\) be any coordinate and let
\(w^{+}=w-\eta_{i}m_{i}e_{i}\) where
\(m_{i}= \beta m_{i}^{\text{old}}+(1-\beta)g_{i}\) and
\(g_{i}= \partial_{i}f(w)\).
Then, using the coordinate‑wise smoothness (Assumption A2) and the
chosen stepsize \(\eta_{i}=1/(L_{i}(1-\beta))\),
\[
f(w^{+})\le f(w)-\frac{1-\beta}{2L_{i}}\,g_{i}^{2}
          -\frac{\beta(1-\beta)}{2L_{i}}\,\bigl(g_{i}^{\text{old}}\bigr)^{2},
\tag{L1}
\]
where \(g_{i}^{\text{old}}=\partial_{i}f(w^{\text{old}})\) is the gradient used
in the previous momentum update.
\end{lemma}
\begin{proof}
From A2,
\[
f\bigl(w-\eta_{i}m_{i}e_{i}\bigr)
\le f(w)-\eta_{i}g_{i}m_{i}+\frac{L_{i}}{2}\eta_{i}^{2}m_{i}^{2}.
\tag{P1}
\]
Insert \(m_{i}= \beta m_{i}^{\text{old}}+(1-\beta)g_{i}\) and
\(\eta_{i}=1/(L_{i}(1-\beta))\).  After expanding the quadratic term and
collecting the coefficients of \(g_{i}^{2}\) and
\((g_{i}^{\text{old}})^{2}\) one obtains exactly the right‑hand side of
(L1).  The algebraic manipulations are displayed step‑by‑step in the
main proof below.  ∎
\end{proof}

-----------------------------------------------------------------------

\section*{Main Proof (Step‑by‑Step Derivation)}
We prove Theorem~\ref{thm:linear}.  All expectations are taken with
respect to the random index \(i_{k}\) conditioned on the past
\(\mathcal{F}_{k}=\sigma\{i_{0},\dots,i_{k-1}\}\).

\noindent\textbf{Notation.}
\(w^{k}\) – iterate after \(k\) updates,  
\(g_{i}^{k}= \partial_{i}f(w^{k})\),  
\(m_{i}^{k}\) – momentum of coordinate \(i\) after iteration \(k\),  
\(e_{i}\) – \(i\)-th unit vector,  
\(w^{*}\) – unique minimiser of \(f\).

\begin{enumerate}
\item[Step 1.] \textbf{Apply Lemma~\ref{lem:one-step} conditioned on the
sampled index.}  Using (L1) with \(i=i_{k}\) and taking conditional
expectation,
\[
\E\!\left[f(w^{k+1})\mid\mathcal{F}_{k}\right]
\le f(w^{k})-
\sum_{i=1}^{d}p_{i}\frac{1-\beta}{2L_{i}}\,\bigl(g_{i}^{k}\bigr)^{2}
-
\sum_{i=1}^{d}p_{i}\frac{\beta(1-\beta)}{2L_{i}}\,
      \bigl(g_{i}^{k-1}\bigr)^{2}.
\tag{S1}
\]

\item[Step 2.] \textbf{Relate the squared gradients to function suboptimality.}
Strong convexity (A1) implies
\[
\frac{1}{2L_{i}}\bigl(g_{i}^{k}\bigr)^{2}
\;\ge\;
\frac{\mu}{L_{i}}\bigl(f(w^{k})-f(w^{*})\bigr)
\quad\text{for each }i,
\tag{S2}
\]
which follows from the inequality
\(\|\nabla f(w^{k})\|^{2}\ge 2\mu\bigl(f(w^{k})-f(w^{*})\bigr)\) and the
coordinate‑wise bound \(|g_{i}^{k}|\le \|\nabla f(w^{k})\|\).

\item[Step 3.] \textbf{Insert (S2) into (S1).}
Using \(\kappa_{i}=L_{i}/\mu\) and the definition
\(\kappa_{\max}=\max_{i}\kappa_{i}\),
\[
\E\!\left[f(w^{k+1})\mid\mathcal{F}_{k}\right]
\le f(w^{k})-
(1-\beta)\,\underbrace{\Bigl(\min_{i}p_{i}\frac{1}{\kappa_{i}}\Bigr)}_{\displaystyle
\ge\,\frac{1}{\kappa_{\max}}\min_{i}p_{i}}\,
\bigl(f(w^{k})-f(w^{*})\bigr)
-
\beta(1-\beta)\,\underbrace{\Bigl(\min_{i}p_{i}\frac{1}{\kappa_{i}}\Bigr)}_{\ge
\frac{1}{\kappa_{\max}}\min_{i}p_{i}}\,
\bigl(f(w^{k-1})-f(w^{*})\bigr).
\tag{S3}
\]

\item[Step 4.] \textbf{Define the contraction factor.}
Set
\[
\rho\;:=\;\frac{1-\beta}{\kappa_{\max}}\;\min_{i}p_{i}\;\in\;(0,1).
\tag{S4}
\]
Then (S3) becomes
\[
\E\!\left[f(w^{k+1})\mid\mathcal{F}_{k}\right]
\le f(w^{k})-\rho\bigl(f(w^{k})-f(w^{*})\bigr)
      -\beta\rho\bigl(f(w^{k-1})-f(w^{*})\bigr).
\tag{S5}
\]

\item[Step 5.] \textbf{Take total expectation and drop the non‑negative term.}
Since the last term in (S5) is non‑negative,
\[
\E\!\left[f(w^{k+1})-f(w^{*})\right]
\le (1-\rho)\,\E\!\left[f(w^{k})-f(w^{*})\right].
\tag{S6}
\]

\item[Step 6.] \textbf{Unroll the recursion.}
Iterating (S6) from \(0\) to \(k-1\) yields
\[
\E\!\left[f(w^{k})-f(w^{*})\right]
\le (1-\rho)^{k}\,\bigl(f(w^{0})-f(w^{*})\bigr).
\tag{S7}
\]

\item[Step 7.] \textbf{Convert function error to distance.}
Strong convexity (A1) gives
\[
\frac{\mu}{2}\,\|w^{k}-w^{*}\|^{2}
\le f(w^{k})-f(w^{*}),
\]
hence
\[
\E\!\bigl[\|w^{k}-w^{*}\|^{2}\bigr]
\le \frac{2}{\mu}\,(1-\rho)^{k}\,
      \bigl(f(w^{0})-f(w^{*})\bigr).
\tag{S8}
\]

\end{enumerate}
Equations (S7) and (S8) are exactly the statements of
Theorem~\ref{thm:linear}.  ∎

-----------------------------------------------------------------------

\section*{Rate Derivation (With Constants)}
The contraction factor derived in \eqref{S4} can be written explicitly as
\[
\rho
=
\frac{1-\beta}{\displaystyle\max_{i}\frac{L_{i}}{\mu}}
\;\min_{i}p_{i}
=
\frac{(1-\beta)\,\mu}{L_{\max}}\;\min_{i}p_{i},
\qquad
L_{\max}:=\max_{i}L_{i}.
\]
Thus the expected suboptimality after \(k\) iterations satisfies
\[
\E\!\left[f(w^{k})-f(w^{*})\right]
\le
\Bigl(1-\frac{(1-\beta)\,\mu}{L_{\max}}\min_{i}p_{i}\Bigr)^{k}
\bigl(f(w^{0})-f(w^{*})\bigr).
\]
If the probabilities are chosen proportionally to the inverse smoothness,
\(p_{i}=L_{\max}/(dL_{i})\), then
\(\min_{i}p_{i}=L_{\max}/(dL_{\max})=1/d\) and
\[
\rho=\frac{1-\beta}{d}\,\frac{\mu}{L_{\max}}.
\]
Hence the method attains a linear rate that improves by a factor
\(1/(1-\beta)\) over classical stochastic coordinate descent.

-----------------------------------------------------------------------

\section*{Special Cases}
\begin{itemize}
\item \textbf{No momentum (\(\beta=0\)).}  The algorithm reduces to the
standard stochastic coordinate descent with stepsize
\(\eta_{i}=1/L_{i}\).  Theorem~\ref{thm:linear} recovers the well‑known
rate \(\bigl(1-\frac{\mu}{L_{\max}}\min_{i}p_{i}\bigr)^{k}\).

\item \textbf{Uniform sampling and identical smoothness.}
  If \(p_{i}=1/d\) and \(L_{i}=L\) for all \(i\),
  \[
  \rho=\frac{1-\beta}{d}\frac{\mu}{L},
  \]
  i.e. each iteration contracts the error by a factor
  \(1-\frac{1-\beta}{d\kappa}\) where \(\kappa=L/\mu\) is the global
  condition number.

\item \textbf{Importance sampling.}
  Choosing \(p_{i}\propto 1/L_{i}\) maximises \(\min_{i}p_{i}\) and
  consequently maximises \(\rho\).  In the extreme case where the
  probabilities are exactly \(p_{i}=L_{\max}/(dL_{i})\) the bound becomes
  \(\rho=(1-\beta)\mu/(dL_{\max})\).

\end{itemize}

-----------------------------------------------------------------------

\end{document}